\documentclass[11pt]{article}
\usepackage[margin=1in]{geometry}
\usepackage{booktabs}
\usepackage{array}
\usepackage{hyperref}
\usepackage{caption}
\usepackage{enumitem}
\usepackage{times}

\title{Milestone 2 Progress Report\\Robust AI-Generated Image Detection under Cross-Generator and Cross-Dataset Distribution Shift}
\author{}
\date{}

\begin{document}
\maketitle

\section{Project Goal and Current Progress}

This project studies binary image classification under distribution shift: each image is classified as either real (0) or AI-generated (1). The goal is not only to obtain high in-distribution accuracy, but also to measure how well baseline detectors transfer across generators, datasets, and common post-processing transformations.

By Milestone 2, the project has implemented baseline models, established repeatable data preparation, training, evaluation, checkpointing, prediction logging, metric reporting, and plotting procedures, and obtained initial results on CIFAKE, a GenImage subset, robustness transforms, and a WildFake subset. The completed results are summarized in \texttt{outputs/metrics/summary\_all\_experiments.csv} and \texttt{docs/RESULTS\_SUMMARY.md}.

Two caveats are important. First, the GenImage experiment is a fallback available-generator split, not the official default Split B, because \texttt{Stable Diffusion V1.4} is missing locally. Second, the WildFake evaluation is not full WildFake; it is a subset consisting of real images from CelebA-HQ and fake images generated by DDIM.

\section{Implemented Baseline Models}

Two baseline model families have been implemented. The first is a ResNet50 image classifier trained end-to-end for binary real/fake prediction. This serves as a conventional convolutional baseline and was used for the CIFAKE sanity check, the GenImage fallback cross-generator experiment, and the WildFake subset external evaluation.

The second baseline is a frozen CLIP ViT-B/32 encoder with an MLP classification head. In this model, CLIP image features are extracted from a frozen pretrained encoder, cached under \texttt{outputs/features/}, and then used by a lightweight MLP classifier. This model tests whether pretrained vision-language representations improve generator transfer and robustness relative to ResNet50.

\section{Training and Evaluation Procedures}

The project uses CSV-based inventories and splits. Each inventory row stores image path, label, dataset, split, generator/source metadata, image dimensions, format, and validity status. Data splits and evaluation outputs are stored under \texttt{outputs/splits/}, \texttt{outputs/metrics/}, \texttt{outputs/predictions/}, \texttt{outputs/figures/}, \texttt{outputs/checkpoints/}, and \texttt{outputs/logs/}.

CIFAKE was used as a full sanity check with 50,000 real and 50,000 fake training images, plus 10,000 real and 10,000 fake test images. For GenImage, the available local subset contains 58,296 inventory rows: 40,796 real and 17,500 fake images. Available generators are ADM, BigGAN, GLIDE, Midjourney, Stable Diffusion V1.5, VQDM, and Wukong. Since Stable Diffusion V1.4 is missing, a fallback split was used: training on Stable Diffusion V1.5, ADM, and BigGAN, and testing on unseen generators GLIDE, Midjourney, Wukong, and VQDM.

Robustness evaluation was performed on the GenImage fallback unseen-generator test split using JPEG compression, resize down-up, and Gaussian blur. External transfer was evaluated on a WildFake subset with 95,713 valid images: 30,000 real CelebA-HQ images and 65,713 fake DDIM images. The main WildFake subset evaluation split contains 4,000 images balanced as 2,000 real and 2,000 fake.

\section{Initial Quantitative Results}

\begin{table}[h]
\centering
\caption{Main completed experiment results.}
\small
\begin{tabular}{llrrrrr}
\toprule
Experiment & Model & Acc. & Prec. & Recall & F1 & AUROC \\
\midrule
CIFAKE sanity & ResNet50 & 0.973750 & 0.964506 & 0.983700 & 0.974009 & 0.997229 \\
GenImage fallback & ResNet50 & 0.905000 & 0.983294 & 0.824000 & 0.896627 & 0.978279 \\
GenImage fallback & CLIP-MLP & 0.933000 & 0.965091 & 0.898500 & 0.930606 & 0.983852 \\
WildFake subset & ResNet50 & 0.471250 & 0.485199 & 0.942500 & 0.640612 & 0.254992 \\
WildFake subset & CLIP-MLP & 0.504250 & 0.502177 & 0.980500 & 0.664183 & 0.744117 \\
\bottomrule
\end{tabular}
\end{table}

The CIFAKE result confirms that the baseline pipeline can learn a strong real/fake classifier in a controlled setting. On the GenImage fallback unseen-generator split, both models generalize reasonably well, and CLIP-MLP performs best overall, improving F1 from 0.896627 to 0.930606 and AUROC from 0.978279 to 0.983852 relative to ResNet50.

The WildFake subset external evaluation reveals a much stronger cross-dataset distribution shift. ResNet50 collapses on this subset, with AUROC below 0.5. CLIP-MLP also drops substantially compared with GenImage, but remains stronger than ResNet50 with AUROC 0.744117.

\section{Robustness and External Dataset Analysis}

\begin{table}[h]
\centering
\caption{Robustness results on the GenImage fallback unseen-generator test split.}
\small
\begin{tabular}{llrrr}
\toprule
Model & Transform & Clean AUROC & Processed AUROC & Drop \\
\midrule
ResNet50 & clean & 0.978279 & 0.978279 & 0.000000 \\
ResNet50 & JPEG & 0.978279 & 0.936216 & 0.042062 \\
ResNet50 & resize & 0.978279 & 0.955641 & 0.022637 \\
ResNet50 & blur & 0.978279 & 0.958654 & 0.019625 \\
CLIP-MLP & clean & 0.983852 & 0.983852 & 0.000000 \\
CLIP-MLP & JPEG & 0.983852 & 0.968515 & 0.015337 \\
CLIP-MLP & resize & 0.983852 & 0.960458 & 0.023394 \\
CLIP-MLP & blur & 0.983852 & 0.959245 & 0.024608 \\
\bottomrule
\end{tabular}
\end{table}

Both models are affected by post-processing. ResNet50 is most sensitive to JPEG compression, with an AUROC drop of 0.042062. CLIP-MLP has smaller degradation under JPEG compression but is most affected by Gaussian blur among the tested transforms. Overall, CLIP-MLP is more robust across these perturbations. The external WildFake subset result is the clearest evidence of distribution shift; it differs in generator family, real-image source, and likely image statistics, so it is a useful early stress test but not a full WildFake benchmark.

\section{Limitations and Future Work}

The main limitation is incomplete dataset coverage. The official GenImage default Split B has not been run because Stable Diffusion V1.4 is missing, so the current GenImage results are from a fallback available-generator split. Similarly, WildFake was evaluated only on a CelebA-HQ/DDIM subset, not on the full WildFake dataset. The experiments also use a single random seed and do not include a full hyperparameter search.

Future work should first download the missing Stable Diffusion V1.4 subset and run the official GenImage Split B. The WildFake evaluation should be expanded beyond the CelebA-HQ/DDIM subset. Model robustness could be improved with stronger augmentation, especially JPEG compression and blur during training. Additional seeds should be run to measure variance, and error case visualization and calibration analysis should be added to better understand failure modes and improve cross-dataset generalization.

\end{document}
